\clearpage
\fancyhead[L]{\small COMET.}
\SourceLeaf{422}
\begin{center}
\Large COMET.

\medskip
Göttingische gelehrte Anzeigen. No. 28. pp. 273..278. 18 February 1819.
\end{center}

\par\indent We are particularly glad to report on the two comets that \textsc{Pons} discovered in Marseille in November of last year, for they are among the most remarkable in several respects. One of them appears to have been seen first in Germany at the Mannheim observatory. Professor \textsc{Nicolai} made the following observations there and communicated them to Councillor \textsc{Gauss} in a letter of 8 January:

\begin{center}
\setlength{\tabcolsep}{2.2pt}\begin{tabular}{c|ccc|ccc|ccc}
1818 & \multicolumn{3}{c|}{\shortstack{Mean time\\in Mannheim}} & \multicolumn{3}{c|}{\shortstack{Right\\ascension}} & \multicolumn{3}{c}{\shortstack{Northern\\declination}} \\
\hline
Dec. 22 & $6^{\mathrm u}$ & $51^{\mathrm m}$ & $48^{\mathrm s}$ & $326^\circ$ & $18'$ & $13''$ & $2^\circ$ & $54'$ & $5''$ \\
        & 8     & 45     & 37     & 326         & 16    & 59     & 2         & 52    & 24    \\
     23 & 7     & 11     & 25     & 326         & 3     & 39     & 2         & 40    & 26    \\
     24 & 7     & 14     & 3      & 325         & 48    & 33     & 2         & 26    & 59    \\
     25 & 6     & 48     & 49     & 325         & 32    & 49     & 2         & 13    & 34    \\
     29 & 7     & 0      & 31     & 324         & 19    & 58     & 1         & 10    & 57    \\
\end{tabular}
\end{center}

\Needspace{10\baselineskip}\par\indent The following parabolic elements were enclosed with the letter. Professor \textsc{Nicolai} had calculated them provisionally from the observations of 22, 25, and 29 December.

\begin{center}
\begin{tabular}{l l}
Time of perihelion, Jan. 1819 . . . . . . . & 25.0000 in M. \\
Logarithm of the least distance . . . . . . & 0.52933 \\
Longitude of perihelion . . . . . . . . . . & $145^\circ$ $49'$ 10 \\
Longitude of the ascending node . . . . . . . & 330 \quad 14 \quad 17 \\
Inclination of the orbit . . . . . . . . . . . . & 14 \quad 59 \quad 6 \\
\multicolumn{2}{l}{Direction: \emph{prograde.}} \\
\end{tabular}
\end{center}

\Needspace{14\baselineskip}\par\indent Professor \textsc{Harding} made the following observations at our observatory:

\begin{center}
\setlength{\tabcolsep}{2.2pt}\begin{tabular}{c|ccc|ccc|ccc l}
& \multicolumn{3}{c|}{\shortstack{Mean time\\in Göttingen}} & \multicolumn{3}{c|}{\shortstack{Right\\ascension}} & \multicolumn{3}{c}{Declination} & \\
\hline
1818 Dec. 25 & $6^{\mathrm u}$ & $39^{\mathrm m}$ & $52^{\mathrm s}$ & $325^\circ$ & $33'$ & $30''$ & $2^\circ$ & $14'$ & $49''$ & North \\
          26 & 5     & 58     & 32     & 325         & 17    & 8      & 2         & 0     & 31     & \\
          27 & 6     & 7      & 10     & 324         & 59    & 15     & 1         & 45    & 5      & \\
          28 & 6     & 24     & 8      & 324         & 40    & 17     & 1         & 27    & 31     & \\
1819 Jan.  1 & 6     & 36     & 16     & 323         & 11    & 48     & 0         & 14    & 58     & \\
           7 & 6     & 39     & 41     & 319         & 53    & 59     & 2         & 19    & 59     & South \\
           8 & 7     & 37     & 34     & 319         & 15    & 28     & 2         & 51    & 49     & \\
          12 & 6     & 6      & 7      & 315         & 36    & 53     & 5         & 35    & 24     & \\
\end{tabular}
\end{center}

\Needspace{11\baselineskip}\par\indent Professor \textsc{Encke} observed this comet five times at the Seeberg observatory:

\begin{center}
\setlength{\tabcolsep}{2.2pt}\begin{tabular}{c|ccc|ccc|ccc l}
& \multicolumn{3}{c|}{\shortstack{Mean time\\in Seeberg}} & \multicolumn{3}{c|}{\shortstack{Right\\ascension}} & \multicolumn{3}{c}{Declination} & \\
\hline
1819 Jan. 1 & $6^{\mathrm u}$ & $38^{\mathrm m}$ & $46^{\mathrm s}$ & $323^\circ$ & $11'$ & $36''$ & $0^\circ$ & $14'$ & $38''$ & North \\
          4 & 7     & 26     & 56     & 321         & 43    & 43     & 0         & 54    & 59     & South \\
          5 & 7     & 12     & 12     & 321         & 10    & 3      & 1         & 20    & 57     & \\
          6 & 5     & 49     & 48     & 320         & 35    & 14     & 1         & 47    & 57     & \\
         12 & 6     & 16     & 32     & 315         & 35    & 20     & 5         & 35    & 26     & \\
\end{tabular}
\end{center}

\SourceLeaf{423}
\expandafter\def\csname HistoricalNoteText423\endcsname{probably Seeberg time, though Mr. \textsc{Encke} did not state this explicitly.}
\par\indent The results that Professor \textsc{Encke} obtained from his own observations and those made at Mannheim, and reported to Councillor \textsc{Gauss} in a letter of 5 February, deserve the closest attention. Although the observations cover only three weeks, they could not be represented by a parabola without errors of several minutes, which were inadmissible for these observations. After several attempts, Mr. \textsc{Encke} took into account the admittedly very rough data from \textsc{Pons} and found an elliptical orbit with a period of 3 years and 7 months that agreed remarkably well with all the observations.

\Needspace{16\baselineskip}\par\indent The following table shows how the positions observed at Mannheim and Seeberg differ from the calculated positions:

\begin{center}
\begin{tabular}{c|c|c|l}
& Right asc. & Declination & Observer \\
\hline
Dec. 22 & $+\ 7''6$  & $-\ 4''7$  & Nicolai \\
     23 & $+\ 0.3$   & $+17.3$    & --- \\
     24 & $-\ 1.0$   & $+11.7$    & --- \\
     25 & $+11.2$    & $-\ 5.5$   & --- \\
     29 & $+14.6$    & $-\ 1.8$   & --- \\
Jan.  1 & $+15.6$    & $+14.8$    & Encke \\
      4 & $+26.5$    & $+27.2$    & --- \\
      5 & $+30.6$    & $+12.9$    & --- \\
      6 & $+36.1$    & $+23.8$    & --- \\
     12 & $+\ 7.1$   & $+\ 3.0$   & --- \\
\end{tabular}
\end{center}

\par\indent An appreciably better agreement could have been obtained by leaving out \textsc{Pons}’s data.

\Needspace{14\baselineskip}\par\indent The elliptical elements are as follows:

\begin{center}
\begin{tabular}{ll}
Time of perihelion, January 1819. . . & $27.13417$\footnotemark \\
Longitude of perihelion . . . . . & $156^\circ\;14'\;8''$ \\
Longitude of the node . . . . . . . & $334\;18\;8$ \\
\multicolumn{2}{l}{(both referred to the mean equinox of 1819)} \\
\end{tabular}\par\smallskip
\begin{tabular}{ll}
Inclination of the orbit . . . . . . . & $13^\circ\;42'\;30''$ \\
Logarithm of the semimajor axis . & $0.3697758$ \\
Eccentricity angle . . . . . . & $58^\circ\;57'\;24''$ \\
\end{tabular}
\end{center}
\footnotetext{\csname HistoricalNoteText423\endcsname}

\par\indent We hope that earlier \emph{good} observations from France or Italy will come to light, so that the uncertainty of the orbit determination can be narrowed. It would be more important still if some further good observations could be made in southern regions after the passage through perihelion. There is no hope of that at our latitudes, where the comet will rise only a little above the southern horizon.

\Needspace{15\baselineskip}\par\indent The close agreement between these elements and those of the first comet of 1805 is most remarkable. For comparison, we give the elements as Professor \textsc{Bessel} then calculated them for a parabolic orbit. We add only that even at the time, comparison of the observations with the calculation showed almost unmistakable traces of a departure of that comet’s orbit from a parabola, which no astronomer seems to have investigated further.

\begin{center}
\begin{tabular}{l l}
\multicolumn{2}{l}{Passage through perihelion, 1818 November} \\
\multicolumn{2}{l}{18.13782 Paris time} \\
Longitude of perihelion . . . . . . . . . . . . & $147^\circ$ $51'$ $28''$ \\
Longitude of the ascending comet . . . . . . . . & 344\ \ 37\ \ 19 \\
Inclination of the orbit . . . . . . . . . . . . . . & 15\ \ 36\ \ 36 \\
Logarithm of the least distance . . . . . . . & 9.57820 \\
\end{tabular}
\end{center}

\SourceLeaf{424}\noindent (With \textsc{Encke}’s elements for the comet of 1818. 1819, this logarithm comes to 9.52579). It certainly deserves careful investigation whether the present comet may have undergone perturbations since 1805 sufficient to explain the differences in the elements, and whether the two comets should accordingly be considered identical.

\par\indent Professor \textsc{Bessel} independently found the other comet discovered by \textsc{Pons} at the Königsberg observatory on 22 December.

\Needspace{17\baselineskip}\par\indent In a letter of 14 January, this astronomer reported the following observations to Councillor \textsc{Gauss}:

\begin{center}
\setlength{\tabcolsep}{2.2pt}\begin{tabular}{c|ccc|ccc|ccc}
& \multicolumn{3}{c|}{\shortstack{Mean time\\in Königsberg}} & \multicolumn{3}{c|}{\shortstack{Right\\ascension}} & \multicolumn{3}{c}{\shortstack{Northern\\declination}} \\
\hline
1818 Dec. 22 & $7^{\mathrm u}$ & $6^{\mathrm m}$  & $51^{\mathrm s}$ & $303^\circ$ & $1'$  & $21''7$ & $36^\circ$ & $48'$ & $20''2$ \\
             & 7     & 53     & 35     & 303         & 10    & 14.7    & 36         & 48    & 29.8    \\
             & 10    & 21     & 47     & 303         & 37    & 29.1    & 36         & 51    & $0.3\!:$ \\
          24 & 18    & 22     & 50     & 311         & 56    & $29.1\!:$ & \multicolumn{3}{c}{.\quad .\quad .\quad .} \\
          25 & 6     & 22     & 20     & 313         & 17    & 17.2    & 37         & 7     & 53.1    \\
          26 & 6     & 4      & 56     & 315         & 38    & 48.2    & 37         & 4     & 1.2     \\
             & 10    & 33     & 0      & 316         & 3     & 18.1    & 37         & 2     & 35.2    \\
          27 & 6     & 55     & 0      & 317         & 39    & 48.8    & 36         & 58    & 45.6    \\
          28 & 6     & 0      & 22     & 319         & 24    & $50.7\!:$ & 36       & 53    & $31.6\!:$ \\
1819 Jan.  1 & 11    & 11     & 4      & 324         & 39    & 12.3    & 36         & 22    & $34.2\!::$ \\
           2 & 6     & 5      & 6      & 325         & 22    & 3.4     & 36         & 15    & 54.3    \\
\end{tabular}
\end{center}

\Needspace{12\baselineskip}\par\indent The following provisional parabolic elements, which Professor \textsc{Bessel} based on the observations of 22 and 27 December and 2 January, were enclosed with the letter:

\begin{center}
\begin{tabular}{l l}
\multicolumn{2}{l}{Time of perihelion, 1818 December 4.0968 Paris time.} \\
Ascending node . . . . . . . . . . . . & $90^\circ$ $7'$ $29''$ \\
Inclination of the orbit . . . . . . . . . . . . . & 117\ \ 19\ \ 10 \\
Distance of perihelion from the node . . . . . . . . & 347\ \ \ 0\ \ 24 \\
Logarithm of the least distance . . . . . . . & 9.928324 \\
\end{tabular}
\end{center}

\Needspace{14\baselineskip}\par\indent Guided by these elements, Councillor \textsc{Gauss} immediately found this comet as a very faint nebulous patch. The elements gave a right ascension too large by about 40 and a declination too large by $10'$. Professor \textsc{Harding} then made the following further observations at our observatory:

\begin{center}
\setlength{\tabcolsep}{2.2pt}\begin{tabular}{c|ccc|ccc|ccc}
& \multicolumn{3}{c|}{\shortstack{Mean time\\in Göttingen}} & \multicolumn{3}{c|}{\shortstack{Right\\ascension}} & \multicolumn{3}{c}{\shortstack{Northern\\declination}} \\
\hline
1819 Jan. 26 & $7^{\mathrm u}$ & $19^{\mathrm m}$ & $5^{\mathrm s}$  & $335^\circ$ & $23'$ & $7''$  & $35^\circ$ & $18'$ & $29''$ \\
          27 & 7     & 2      & 35     & 335         & 36    & 4      & 35         & 18    & 50     \\
          28 & 7     & 22     & 16     & 335         & 50    & 38     & 35         & 20    & 58     \\
          29 & 7     & 27     & 56     & 336         & 4     & 38     & 35         & 22    & 32     \\
          30 & 6     & 54     & 10     & 336         & 16    & 33     & 35         & 23    & 31     \\
\end{tabular}
\end{center}

\par\indent This comet, remarkable for its large geocentric motion since its first discovery, is growing ever fainter, so further observations can hardly be expected. A future investigation will show whether the observations available for this comet too are sufficient to determine the departure of its orbit from a parabola, and what the results will be.
